\pdfbookmark[1]{Abstract}{Abstract}
\thispagestyle{empty}
\begin{center}
  \rule{\textwidth}{0.4pt}\\[0.9em]
  {\Large\spacedallcaps{Abstract}}\\[0.7em]
  \rule{\textwidth}{0.4pt}
\end{center}
\bigskip

\noindent Gnascor is the language the orior compiler is written in. It treats every host, a
processor, a database, a web service or a language model, as a black box that answers yes or no
questions, and it reads each answer together with what the answer cost. A question is an
address, a predicate asked at it, and a timeout; the answer is one bit, 1 where the predicate
holds inside the timeout. The timeout is never written by hand. It is derived from the spread
of the host's own costs, measured first with no timeout at all.

Two answers taken side by side make a branch pair, and the pair is a state of a finite state
machine whose output on each transition is a four-letter mnemonic. The mnemonics are the
language's keywords. Each maps to a term the field already has: failover and failback, load
shedding, a circuit breaker, a health probe, the saturation, latency and error signals of
production monitoring, and the superstate of a statechart.

The method has one rule. A predicate holds or it does not, and that answer carries no noise; a
cost is a measurement and always does. The language decides which candidates hold first and
ranks only the survivors by cost, and never adds the two into one score. The ranking reads a
difference as real only where it clears the candidates' own spread.

This research paper states the language, maps each of its parts to the field's term with a
source, collects the experiments that have run on it, and adds three new measurements: how
many bits of a transition each kind of label keeps, how often the derived timeout is beaten on
a real host, and which states repeat from run to run. It then states three hypotheses the
measurements support and the work left open.
\cleardoublepage
